\chapter{Квантовые флуктуации и происхождение структуры}

\lectureoverview{В заключительной лекции ставится вопрос, каким образом вселенная, сглаженная инфляцией, тем не менее может приобрести крошечные неоднородности, из которых вырастают галактики. Ответ строится поэтапно. Сначала универсальные константы определяют планковские масштабы и выявляют огромный разрыв между привычной физикой частиц и квантовой гравитацией. Затем микроволновое небо задаёт цель: вариации температуры и плотности порядка \(10^{-5}\), распределённые по многим масштабам. Нулевые колебания дают неизбежные флуктуации, затухание объясняет, как колебание может замораживаться, а волновое уравнение в расширяющейся вселенной превращает каждую фурье-моду инфлатона именно в такой затухающий осциллятор. Моды замораживаются, когда их физическая длина волны достигает хаббловского масштаба. Когда инфляция заканчивается, эта замороженная вариация приводит к тому, что соседние области разогреваются в немного разное время, превращая микроскопический квантовый шум в первичную картину распределения плотности.}

\section{Какие константы универсальны?}

Мы располагаем необычайно точной квантовой теорией элементарных частиц. Её параметры включают массы кварков и лептонов, константы связи, данные о смешивании, заряды и сам список частиц. При неограниченной вычислительной мощности эта теория описала бы не только эксперименты на ускорителях, но и атомы, молекулы и химию; в принципе её следствия простираются и на биологию, и на физические процессы, лежащие в основе мышления. Однако от этого такие системы не становятся лёгкими для расчёта. Это означает, что Стандартная модель предназначена служить микроскопическим основанием обычного вещества. Гравитация описывается отдельно, на классическом уровне, общей теорией относительности. Обе теории мощны; но ни одна не делает всякое число, возникающее в низкоэнергетической физике, неизбежным.

Отношение массы протона к массе электрона — вполне корректное физическое число, но оно привязано к двум конкретным частицам. Если заменить протон нейтроном или электрон мюоном, получится другое отношение. Подобные параметры могут кодировать глубокую физику, которой мы пока не знаем, но они не демонстрируют своей универсальности. Иначе обстоит дело с тремя размерными константами:
\[
c,\qquad \hbar,\qquad G.
\]
В определяющих их законах вновь и вновь появляется слово \emph{все}. Всё вещество и вся информация подчиняются предельной скорости \(c\). Все квантовые системы подчиняются соотношениям неопределённости, задаваемым \(\hbar\). Все формы энергии гравитируют с одной и той же ньютоновской постоянной \(G\). Именно эта общность и делает эти константы фундаментальными в рамках настоящего рассуждения.

Планк понял важность этой комбинации ещё до того, как стал известен полный смысл квантовой механики. Его постоянная впервые вошла в описание излучения чёрного тела. Уже тогда был виден размерностный факт: из \(c\), \(\hbar\) и \(G\) можно построить единицы длины, времени и массы, не ссылаясь ни на какой конкретный атом или частицу.

\section{Построение планковских масштабов}

Пусть \(L\), \(T\) и \(M\) обозначают размерности длины, времени и массы. Скорость света имеет размерность
\[
[c]=LT^{-1}.
\]
Соотношение неопределённости \(\Delta x\,\Delta p\gtrsim\hbar\) вместе с \(p=mv\) даёт
\[
[\hbar]=ML^2T^{-1}.
\]
Наконец, уравнения Ньютона
\[
F=\frac{Gm_1m_2}{r^2},
\qquad
F=ma
\]
дают
\[
[G]=L^3M^{-1}T^{-2}.
\]
Этот вывод важен, потому что размерностный анализ — не гадание: каждый показатель степени фиксируется требованием сократить нежелательные размерности.

Ищем длину вида
\begin{equation}
\ell=c^\alpha\hbar^\beta G^\gamma.
\label{eq:c10-length-ansatz}
\end{equation}
Показатель степени при массе должен быть нулевым, так что
\[
\beta-\gamma=0.
\]
Показатель степени при времени должен быть нулевым, а показатель при длине должен равняться единице:
\[
-\alpha-\beta-2\gamma=0,
\qquad
\alpha+2\beta+3\gamma=1.
\]
С учётом \(\gamma=\beta\) это превращается в
\[
\alpha+3\beta=0,
\qquad
\alpha+5\beta=1.
\]
Разность этих уравнений даёт \(2\beta=1\), откуда
\[
\beta=\gamma=\frac12,
\qquad
\alpha=-\frac32.
\]
Следовательно, единственная длина, составленная из трёх универсальных констант, равна
\begin{equation}
\ell_{\mathrm P}
=\sqrt{\frac{\hbar G}{c^3}}
\approx 1.6\times10^{-35}\,\mathrm m.
\label{eq:c10-planck-length}
\end{equation}
В лекции сохраняется только порядок величины, \(10^{-35}\,\mathrm m\), однако размерностный вывод фиксирует и саму комбинацию.

Соответствующее время — это время прохождения света:
\begin{equation}
t_{\mathrm P}
=\frac{\ell_{\mathrm P}}{c}
=\sqrt{\frac{\hbar G}{c^5}}
\approx 5.4\times10^{-44}\,\mathrm s.
\label{eq:c10-planck-time}
\end{equation}
С точностью до порядка величины это \(10^{-43}\,\mathrm s\). Оно не является независимой конструкцией, если уже известны планковская длина и универсальная скорость.

Для массы используйте размерностное соотношение \(Et\sim\hbar\) и \(E=mc^2\). При \(t=t_{\mathrm P}\) получаем
\begin{equation}
m_{\mathrm P}
=\frac{\hbar}{c^2t_{\mathrm P}}
=\sqrt{\frac{\hbar c}{G}}
\approx 2.2\times10^{-8}\,\mathrm{kg}.
\label{eq:c10-planck-mass}
\end{equation}
Это по меркам физики частиц удивительно макроскопическая величина, примерно масса мелкой пылинки. Её энергия покоя имеет порядок \(10^9\,\mathrm J\), сравнимый по масштабу с химической энергией в баке топлива. Автомобиль может проехать сотни миль всего на нескольких баках топлива, потому что при сгорании высвобождается лишь ничтожная доля энергии покоя топлива. Полное превращение одной планковской массы в энергию было бы совсем другим делом. Крошечная масса может соответствовать огромной энергии, когда её умножают на \(c^2\).

Контраст с ускорительной физикой разителен. Расстояние порядка \(10^{-18}\,\mathrm m\), характерное для кратчайших масштабов, исследованных в столкновениях высоких энергий на момент лекции, всё ещё составляет около \(10^{17}\) планковских длин. Поэтому неизведанная область между электрослабым и планковским масштабами огромна. Ядерная физика даёт правильную аналогию. Ядро издалека выглядит элементарным, но на меньших расстояниях оно разрешается на протоны и нейтроны, а те, в свою очередь, — на кварки и глюоны. Низкоэнергетические частицы могут быть простыми по сравнению с ядрами и всё же оставаться сложными эмерджентными объектами по сравнению с тем, что бы ни находилось на планковском масштабе. Многочисленные параметры Стандартной модели выглядят менее удивительными, если смотреть через эту иерархию, хотя их наблюдаемые значения и тонкие подстройки по-прежнему не объяснены.

\begin{classroomqa}
\audiencequestion{Бозон Хиггса легче или тяжелее ожидаемого, и почему его масса кажется загадочной?}
\lecturerresponse{Здесь есть два сравнения. По отношению к планковской массе бозон Хиггса с массой \(125\,\mathrm{GeV}\) необычайно лёгок. В теории без защитной структуры квантовые поправки обычно делают массу скаляра чувствительной к гораздо более высоким масштабам, так что получение наблюдаемого значения выглядит как балансировка параметров на лезвии ножа. По отношению к некоторым простым суперсимметричным ожиданиям, обсуждавшимся тогда, наблюдаемый бозон Хиггса, напротив, можно было бы назвать несколько тяжёлым. Кажущееся противоречие возникает из-за сравнения его с двумя разными шкалами отсчёта; лежащая в основе проблема иерархии остаётся. \lecturetimestamp{00:26:45--00:30:54}}
\end{classroomqa}

\begin{classroomqa}
\audiencequestion{Если Хиггс помогает придавать частицам массу, должны ли частицы тёмной материи иметь массы, близкие к планковскому масштабу?}
\lecturerresponse{Нет. WIMP — это \emph{слабо взаимодействующая массивная частица}; само выражение указывает на такой масштаб взаимодействия, который мог бы сделать возможным обнаружение в лаборатории, а не на планковскую массу. При обычных предположениях о WIMP массы могут быть намного больше массы протона и всё же фантастически меньше \(m_{\mathrm P}\). Если допустить, что взаимодействие с обычным веществом произвольно слабее, космология допускает гораздо более широкий диапазон масс, но прямое обнаружение соответственно становится труднее. Протон и нейтрон тоже дают полезное предостережение: большая часть их массы — это энергия связи сильного взаимодействия, а не просто сумма порождённых Хиггсом масс кварков. \lecturetimestamp{00:30:55--00:33:11}}
\end{classroomqa}

\begin{classroomqa}
\audiencequestion{Что означает планковская масса, и существует ли фундаментальная энергия на бит?}
\lecturerresponse{Полезная гравитационная интерпретация состоит в том, что чёрную дыру нельзя сделать параметрически меньше планковского масштаба, если она должна оставаться обычной полуклассической чёрной дырой. Такой объект имеет массу порядка \(m_{\mathrm P}\), площадь горизонта порядка одной планковской площади и энтропию порядка единицы. В этом ограниченном смысле, относящемся к чёрным дырам, можно сопоставить энергию планковского порядка информации порядка единицы. Это не универсальное утверждение о том, что каждый хранимый бит в любом устройстве обходится в планковскую энергию; обычные частицы и обычные устройства памяти могут быть намного легче. \lecturetimestamp{00:33:11--00:35:10}}
\end{classroomqa}

Начиная с этого места удобно использовать планковские единицы,
\begin{equation}
c=\hbar=G=1.
\label{eq:c10-planck-units}
\end{equation}
Тогда длины, времена, массы и энергии можно выражать чистыми числами, измеренными относительно формул~\eqref{eq:c10-planck-length}--\eqref{eq:c10-planck-mass}.

\section{Почти однородное микроволновое небо}

Космологическая проблема начинается с микроволновой карты всего неба. WMAP, Wilkinson Microwave Anisotropy Probe, измерил излучение, высвободившееся вблизи эпохи, когда первичная плазма стала прозрачной. С тех пор фотоны испытали красное смещение примерно в тысячу раз и теперь образуют почти идеальное чёрное тело с температурой около \(2.7\,\mathrm K\). Когда мы смотрим вдаль, мы не видим вселенную в одном общем текущем моменте времени. Чем больше расстояние, тем раньше произошло излучение, а микроволновый фон — это далёкая оболочка, откуда фотоны впервые свободно устремились к нам.

\begin{classroomqa}
\audiencequestion{Что именно изображает овальная карта WMAP?}
\lecturerresponse{Каждая точка в овале соответствует некоторому направлению на небе; сфера здесь сплющена картографической проекцией. Левая и правая границы смыкаются, верхняя и нижняя части представляют противоположные галактические полушария, а галактическая плоскость проходит посередине. Микроволновое излучение переднего плана от нашей собственной Галактики необходимо смоделировать и вычесть. То, что остаётся, — это не карта близких звёзд, а карта микроволновой температуры по направлениям на оболочке последнего рассеяния.\lecturetimestamp{00:37:49--00:40:49}}
\end{classroomqa}

В первом приближении результат изотропен:
\[
T(\hat{\mathbf n})\simeq 2.7\,\mathrm K
\]
для любого направления наблюдения \(\hat{\mathbf n}\). Совершенная изотропия вокруг каждой точки означала бы совершенную однородность. Но совершенно однородная Вселенная не породила бы самопроизвольно галактики, скопления и похожую на каустики космическую паутину, обсуждавшуюся на прошлой лекции. Некоторая неоднородность должна была существовать, как бы мала ни была её амплитуда.

Запишем отклонение температуры в данном направлении как
\[
\delta T(\hat{\mathbf n})=T(\hat{\mathbf n})-\bar T.
\]
Измеренная относительная вариация имеет порядок
\begin{equation}
\frac{\delta T}{T}\sim10^{-5}.
\label{eq:c10-cmb-anisotropy}
\end{equation}
Один путь к этому числу был динамическим: проследить сегодняшнее распределение вещества назад во времени и спросить, какой начальный контраст требуется к моменту расцепления. Другой был прямым: увеличить почти лишённую особенностей микроволновую карту на пять порядков и измерить саму пятнистость. Эти два пути согласуются на уровне, достаточном для данного рассуждения.

Пятнистость наблюдается на многих угловых масштабах. Большие пятна содержат меньшие, и при увеличении области обнаруживается статистически сходная структура. Термин \emph{масштабно-инвариантный} не означает, что каждый масштаб буквально тождествен другому; он означает, что мощность флуктуаций меняется с масштабом лишь слабо в рассматриваемом диапазоне. Эта почти масштабная инвариантность — часть тех данных, которые теория должна объяснить, а не только общая амплитуда.

Почему температура должна отслеживать вещество? Фотон, выходящий из ямы гравитационного потенциала, теряет энергию. Поэтому излучение, приходящее из более плотной области, может быть смещено относительно излучения из менее плотной области. На уровне оценок по порядку величины, используемых в лекции, это мотивирует
\begin{equation}
\frac{\delta\rho}{\rho}\sim10^{-5}
\label{eq:c10-density-target}
\end{equation}
вблизи расцепления. Точный расчёт CMB включает акустическую динамику, гравитационные потенциалы и передаточные функции, так что уравнение~\eqref{eq:c10-density-target} — это краткая запись, задающая масштаб, а не точное равенство между \(\delta T/T\) и \(\delta\rho/\rho\).

\begin{classroomqa}
\audiencequestion{Может ли эта пятнистость быть обычным микроскопическим шумом, возникшим при последнем рассеянии?}
\lecturerresponse{Не на наблюдаемых угловых масштабах. К моменту расцепления эти пятна соответствовали огромным физическим областям, а не возмущениям атомного масштаба в рекомбинирующей плазме. Если, однако, проследить их гораздо дальше назад, их физические длины волн становятся меньше. Инфляция даёт недостающее увеличение: флуктуации, начавшиеся как квантовые микроскопические возмущения, могут растянуться до астрономических масштабов.\lecturetimestamp{00:46:42--00:48:02}}
\end{classroomqa}

Этот механизм можно понять, опираясь на три физических элемента: нулевые колебания гармонического осциллятора, затухание осциллятора и волновое уравнение для скалярного поля в расширяющейся Вселенной. Мы соберём их именно в таком порядке.

\section{Нулевые колебания осциллятора}

Для гармонического осциллятора единичной массы,
\begin{equation}
E=\frac12\dot x^{\,2}+\frac12\omega^2x^2,
\qquad
\ddot x+\omega^2x=0.
\label{eq:c10-undamped}
\end{equation}
Кинетическая и потенциальная энергии непрерывно переходят друг в друга. В среднем за период каждая несёт половину полной энергии.

Квантовая механика даёт уровни энергии
\begin{equation}
E_n=\left(n+\frac12\right)\hbar\omega.
\label{eq:c10-oscillator-levels}
\end{equation}
Основное состояние не может иметь одновременно \(x=0\) и \(\dot x=0\), поскольку это приписало бы одновременно точные координату и импульс. Его неустранимая энергия,
\[
E_0=\frac12\hbar\omega,
\]
есть энергия нулевых колебаний. Слово \emph{колебание} здесь следует понимать статистически: стационарное основное состояние не описывает классическую синусоидальную траекторию, но имеет ненулевые дисперсии координаты и импульса.

В среднем половина энергии основного состояния приходится на потенциальный член:
\[
\frac12\omega^2\langle x^2\rangle
=\frac14\hbar\omega.
\]
Тем самым
\begin{equation}
\langle x^2\rangle=\frac{\hbar}{2\omega}
\sim\frac{\hbar}{\omega}.
\label{eq:c10-zero-point-position}
\end{equation}
Пренебрегая множителями 2, удобная комплексная запись для масштаба флуктуации имеет вид
\[
x(t)\sim\sqrt{\frac{\hbar}{\omega}}\,e^{i\omega t}.
\]
Тогда дифференцирование даёт
\begin{equation}
\langle\dot x^{\,2}\rangle\sim\hbar\omega,
\label{eq:c10-zero-point-velocity}
\end{equation}
в согласии с равенством средних кинетической и потенциальной энергий. Охлаждение может устранить возбуждения с \(n>0\), но не может устранить эти дисперсии основного состояния.

\section{Затухание, критичность и замораживание}

Погрузим осциллятор в вязкую среду. Система пружина—груз, движущаяся в мёде, подчиняется уравнению
\begin{equation}
\ddot x+\gamma\dot x+\omega^2x=0,
\label{eq:c10-damped}
\end{equation}
где \(\gamma>0\) характеризует сопротивление. Ищем решение в виде \(x=e^{\alpha t}\). Тогда получаем
\begin{equation}
\alpha^2+\gamma\alpha+\omega^2=0,
\end{equation}
и, следовательно,
\begin{equation}
\alpha_\pm
=\frac{-\gamma\pm\sqrt{\gamma^2-4\omega^2}}{2}.
\label{eq:c10-damped-roots}
\end{equation}
Эти два корня дают два независимых решения, требуемых для дифференциального уравнения второго порядка.

\begin{classroomqa}
\audiencequestion{Не должен ли член сопротивления в уравнении~\eqref{eq:c10-damped} иметь знак минус?}
\lecturerresponse{Сама сила сопротивления равна \(F_{\rm drag}=-\gamma\dot x\). Уравнение Ньютона имеет вид \(\ddot x=-\gamma\dot x-\omega^2x\). Перенос обоих силовых членов в левую часть даёт уравнение~\eqref{eq:c10-damped} со знаками плюс. Следовательно, знак в дифференциальном уравнении и направление физической силы согласованы.\lecturetimestamp{00:59:38--01:00:10}}
\end{classroomqa}

Если \(\gamma^2<4\omega^2\), корни комплексные. Введём
\[
\Omega=\sqrt{\omega^2-\frac{\gamma^2}{4}}.
\]
Тогда
\begin{equation}
x(t)=e^{-\gamma t/2}
\left(A\cos\Omega t+B\sin\Omega t\right).
\label{eq:c10-underdamped}
\end{equation}
Осциллятор по-прежнему многократно проходит через точку равновесия, но его огибающая затухает. Это режим слабого затухания.

Если \(\gamma^2>4\omega^2\), оба корня в уравнении~\eqref{eq:c10-damped-roots} вещественны и отрицательны:
\begin{equation}
x(t)=A e^{\alpha_+t}+B e^{\alpha_-t}.
\label{eq:c10-overdamped}
\end{equation}
Устойчивых колебаний нет; обе компоненты затухают. Это режим сильного затухания. При
\begin{equation}
\gamma=2\omega
\label{eq:c10-critical}
\end{equation}
корни совпадают, что соответствует критическому затуханию, переходу между двумя типами поведения.

Один предельный случай особенно важен. При отсутствии восстанавливающей силы, \(\omega=0\),
\[
\alpha(\alpha+\gamma)=0,
\]
так что
\begin{equation}
x(t)=c_0+c_1e^{-\gamma t}.
\label{eq:c10-no-restoring}
\end{equation}
Скорость затухает, но координата не обязана возвращаться к нулю. Она замораживается на постоянном значении \(c_0\), определяемом начальными условиями. Камень, протолкнутый через мёд без прикреплённой пружины, останавливается там, где сопротивление полностью гасит его движение.

Теперь позволим частоте медленно меняться. Начнём с \(\omega\gg\gamma\), так что движение колебательное. Уменьшаем \(\omega(t)\), удерживая коэффициент затухания приблизительно постоянным. Колебание замедляется, проходит через \(\gamma=2\omega\), переходит в режим сильного затухания и по мере \(\omega\to0\) приближается к форме без восстанавливающей силы. Эквивалентно можно увеличивать затухание относительно восстанавливающей силы. В космологическом приложении решающим изменением будет падение эффективной частоты, тогда как хаббловское затухание остаётся почти постоянным. Окончательное замороженное значение нельзя вывести только из самого перехода; оно определяется всей предшествующей эволюцией и начальным квантовым состоянием.

\section{Скалярные волны в расширяющейся Вселенной}

Вернёмся к инфлатону \(\phi\). Предполагается, что его потенциал содержит высокое широкое плато, за которым следует крутой спуск к более низкому минимуму. На плато важна аппроксимация
\begin{equation}
V'(\phi)\simeq0,
\label{eq:c10-flat-potential}
\end{equation}
а не \(V(\phi)\simeq0\). Наклон мал, тогда как плотность энергии может оставаться достаточно большой, чтобы обеспечивать инфляцию. Подробное происхождение такой формы не задаётся; это феноменологическое требование модели.

\begin{figure}[htbp]
\centering
\includegraphics[width=0.88\textwidth]{lecture_10_figure_02.png}
\caption{Слева по-прежнему видны плато и обрыв потенциала инфлатона, а справа начинается запись волнового уравнения в плоском пространстве. На доске отчётливо видны приближённое условие \(V'(\phi)=0\) и начало члена со второй производной, но лектор закрывает незавершённое уравнение; уравнение~\eqref{eq:c10-flat-wave} восстановлено по устному выводу.}
\label{fig:c10-potential-wave}
\lectureframe{01:17:17.865}
\end{figure}

Для почти свободного скалярного поля в плоском пространстве-времени,
\begin{equation}
\frac{\partial^2\phi}{\partial t^2}
-\nabla^2\phi=0.
\label{eq:c10-flat-wave}
\end{equation}
Если бы потенциал не был плоским, справа добавился бы член \(-V'(\phi)\). В приближении плато мы его опускаем. Если бы поле было идеально однородным, пространственная производная исчезла бы. В расширяющейся Вселенной оставшееся однородное уравнение имело бы вид
\[
\ddot\phi+3H\dot\phi=0,
\]
это уравнение хаббловского трения из предыдущей лекции.

\begin{figure}[htbp]
\centering
\includegraphics[width=0.88\textwidth]{lecture_10_figure_03.png}
\caption{Слева указано плато, а справа появляются последовательные стадии волнового уравнения. Верхний оператор для плоского пространства читается; нижняя строка доходит до \(+3H\), но не завершена, поэтому полный член затухания взят из сопутствующего устного объяснения.}
\label{fig:c10-hubble-board}
\lectureframe{01:18:57.554}
\end{figure}

\begin{figure}[htbp]
\centering
\begin{coursefitfigure}
\begin{tikzpicture}[x=1cm,y=0.9cm,>=Latex]
  \draw[->] (-0.2,0) -- (9.3,0) node[right] {\(\phi\)};
  \draw[->] (0,-0.2) -- (0,4.6) node[above] {\(V(\phi)\)};
  \draw[very thick,blue!55!black,smooth]
    plot coordinates {(0.25,3.9) (1.4,3.88) (2.8,3.8) (4.2,3.6) (5.3,3.25) (6.0,2.6) (6.5,1.45) (6.9,0.45) (7.5,0.28) (8.5,0.7)};
  \draw[dashed] (6.45,0) -- (6.45,2.0);
  \draw[->,thick,red!65!black] (2.0,4.25) -- (3.1,3.82);
  \node[anchor=south west] at (0.8,4.0) {slow-roll plateau: \(V'\simeq0\)};
  \node[anchor=west] at (6.55,2.05) {end of inflation};
  \node[anchor=north] at (7.45,0.22) {reheating region};
\end{tikzpicture}
\end{coursefitfigure}
\caption{Чистая реконструкция потенциала, используемого на протяжении всего рассуждения. Малый наклон допускает медленное скатывание, тогда как высокое значение \(V\) поддерживает \(H\) почти постоянным; крутой спуск завершает инфляцию и высвобождает запасённую энергию.}
\label{fig:c10-potential-reconstruction}
\end{figure}

Для неоднородного поля необходимо различать координатное и физическое расстояния. Выберем безразмерные сопутствующие метки \(\mathbf x\), так что бесконечно малое собственное расстояние равно
\[
d\boldsymbol\ell=a(t)\,d\mathbf x.
\]
При фиксированном времени,
\[
\nabla_{\rm proper}=\frac{1}{a}\nabla_{\mathbf x},
\qquad
\nabla_{\rm proper}^2=\frac{1}{a^2}\nabla_{\mathbf x}^2.
\]
Следовательно, уравнение для скалярного поля имеет вид
\begin{equation}
\ddot\phi+3H\dot\phi-\frac{1}{a^2}\nabla^2\phi=0,
\label{eq:c10-frw-wave}
\end{equation}
где \(\nabla\) дифференцирует по сопутствующим меткам. Если вместо этого выбрать безразмерный масштабный фактор, то сопутствующие координаты будут иметь размерность длины; физика не меняется.

\begin{classroomqa}
\audiencequestion{Должен ли пространственный член в уравнении~\eqref{eq:c10-frw-wave} быть положительным?}
\lecturerresponse{Нет. Волновой оператор в плоском пространстве уже содержит относительный знак \(\partial_t^2-\nabla^2\). Расширение изменяет пространственную производную множителем \(a^{-2}\), но не обращает этот знак. После подстановки Фурье два знака минус объединяются, и уравнение для моды приобретает ожидаемый положительный восстанавливающий член \(+k^2\phi_k/a^2\). \lecturetimestamp{01:22:18--01:23:06}}
\end{classroomqa}

\section{Каждая мода Фурье — затухающий осциллятор}

Рассмотрим по одной пространственной волне:
\begin{equation}
\phi(\mathbf x,t)
=\phi_{\mathbf k}(t)e^{i\mathbf k\cdot\mathbf x}.
\label{eq:c10-mode-ansatz}
\end{equation}
Комплексная экспонента служит удобной записью для синусов и косинусов. Произвольный вещественный профиль строится суперпозицией мод с соответствующим условием вещественности между \(\mathbf k\) и \(-\mathbf k\).

Производные по времени действуют только на \(\phi_{\mathbf k}(t)\), тогда как
\[
\nabla^2e^{i\mathbf k\cdot\mathbf x}
=-k^2e^{i\mathbf k\cdot\mathbf x}.
\]
Подстановка в уравнение~\eqref{eq:c10-frw-wave} дает
\begin{equation}
\ddot\phi_{\mathbf k}
+3H\dot\phi_{\mathbf k}
+\frac{k^2}{a^2}\phi_{\mathbf k}=0.
\label{eq:c10-mode-equation}
\end{equation}
Таким образом, каждое волновое число ведет себя как затухающий осциллятор с
\begin{equation}
\gamma=3H,
\qquad
\omega_k(t)=\frac{k}{a(t)}.
\label{eq:c10-mode-identification}
\end{equation}
Физическое волновое число \(k/a\) испытывает красное смещение по мере расширения Вселенной. Следовательно, эффективная восстанавливающая сила \(k^2/a^2\) убывает к нулю.

Мода, которая начинается с малой физической длины волны, имеет \(\omega_k\gg H\) и осциллирует. Ее длина волны растягивается, частота падает, и в конце концов начинает доминировать хаббловское затухание. Тогда она стремится к неосциллирующему значению, а не возвращается к нулю. Разные моды \(\mathbf k\) проходят этот переход в разное время: большему \(k\), которому при фиксированном \(a\) соответствует более короткая длина волны, требуется большее расширение, прежде чем его физическая частота упадет до хаббловского масштаба.

\section{Выход за горизонт и замороженные вакуумные флуктуации}

Используя критическое условие осциллятора как оценку порядка единицы,
\[
\gamma=2\omega_k,
\]
уравнение~\eqref{eq:c10-mode-identification} дает
\begin{equation}
3H=\frac{2k}{a},
\qquad
a_{\rm cross}=\frac{2k}{3H}.
\label{eq:c10-critical-mode}
\end{equation}
Суть не в численных множителях. Переход происходит, когда физическое волновое число \(k/a\) имеет порядок \(H\).

Пусть
\[
\lambda_{\rm com}=\frac{2\pi}{k},
\qquad
\lambda_{\rm phys}=a\lambda_{\rm com}
=\frac{2\pi a}{k}.
\]
Проверка размерностей показывает, где должно стоять \(a\): при безразмерном \(\mathbf x\) и \(k\), и \(\mathbf k\cdot\mathbf x\) безразмерны, тогда как \(a\) и \(\lambda_{\rm phys}\) имеют размерность длины. Поэтому из уравнения~\eqref{eq:c10-critical-mode} следует
\begin{equation}
\lambda_{\rm phys}\sim H^{-1},
\label{eq:c10-horizon-crossing}
\end{equation}
с точностью до таких множителей, как \(2\), \(3\) и \(2\pi\).

Закон Хаббла \(v=Hd\) сопоставляет расстоянию \(H^{-1}\) скорость удаления порядка \(c=1\). Во время инфляции \(H\) велик и почти постоянен, так что эта хаббловская длина микроскопична, хотя и играет роль причинного масштаба. То, что это событие называют \emph{выходом за горизонт}, — условность: никакой материальный объект не проходит через физическую стену. Физическая длина волны моды становится больше приблизительно фиксированного хаббловского масштаба, и ее динамика меняется от быстрых осцилляций к почти замороженной амплитуде.

\begin{figure}[htbp]
\centering
\begin{coursefitfigure}
\begin{tikzpicture}[x=1cm,y=0.85cm,>=Latex]
  \draw[->] (0,0) -- (9.2,0) node[right] {cosmic time};
  \draw[->] (0,0) -- (0,5.0) node[above] {physical length};
  \draw[dashed,thick] (0.25,2.75) -- (8.75,2.75)
    node[right] {\(H^{-1}\)};
  \draw[very thick,blue!55!black]
    (0.35,0.55) .. controls (3.3,0.75) and (5.3,1.55) .. (8.5,4.55);
  \node[blue!55!black,anchor=west] at (6.65,4.55) {\(\lambda_{\rm phys}=2\pi a/k\)};
  \draw[densely dotted] (6.05,0) -- (6.05,2.75);
  \fill (6.05,2.75) circle (2pt);
  \node[anchor=north] at (6.05,0) {horizon exit};
  \node[align=center] at (2.6,1.7) {oscillating\\sub-Hubble mode};
\end{tikzpicture}
\end{coursefitfigure}
\caption{При почти экспоненциальном расширении физическая длина волны растет, тогда как хаббловская длина остается приблизительно постоянной. Их пересечение отмечает переход от быстрых осцилляций к замороженной моде с точностью до множителей порядка единицы.}
\label{fig:c10-horizon-reconstruction}
\end{figure}

Вакуумные флуктуации не позволяют инфляции просто стереть каждую волну. Каждая коротковолновая мода начинается с нулевых колебаний. Расширение растягивает одну моду до хаббловского масштаба и замораживает её, тогда как ещё более коротковолновые вакуумные моды продолжают осциллировать. Затем и они растягиваются и замораживаются. Повторение этого процесса создаёт профиль поля, содержащий замороженные вклады в широком диапазоне длин волн и направлений. Результат стохастичен, но не бесструктурен: его статистика определяется квантовым состоянием и инфляционным фоном.

\begin{figure}[htbp]
\centering
\includegraphics[width=0.88\textwidth]{lecture_10_figure_04.png}
\caption{На всей доске показана суперпозиция длинных и коротких волн поля; рядом с метками \(\gamma\) и \(3H\) дан небольшой набросок затухания. Этот рисунок фиксирует качественную картину мод; по нему не считывают точную функцию.}
\label{fig:c10-mode-board}
\lectureframe{01:45:28.476}
\end{figure}

\begin{classroomqa}
\audiencequestion{Замерзают ли моды в точке или при некоторой длине волны, и растёт ли экспоненциально само \(H\)?}
\lecturerresponse{Они замерзают, когда длина волны достигает хаббловского масштаба. На почти плоском плато \(V\) и, следовательно, \(H\) приблизительно постоянны; экспоненциально растёт масштабный фактор \(a(t)\), а значит, и физическая длина волны \(a/k\). Каждая комовинговая мода достигает условия \(k/a\sim H\) в свой момент времени. \lecturetimestamp{01:45:06--01:45:39}}
\end{classroomqa}

\begin{classroomqa}
\audiencequestion{Это электромагнитные волны?}
\lecturerresponse{Нет. Это флуктуации скалярного поля инфлатона. Электромагнитные вакуумные флуктуации служат привычной аналогией, но переменная в уравнении~\eqref{eq:c10-mode-equation} — это \(\phi\). \lecturetimestamp{01:45:39--01:46:04}}
\end{classroomqa}

\begin{classroomqa}
\audiencequestion{Существует ли континуум длин волн?}
\lecturerresponse{В идеализации с бесконечным пространством — да. Поле представляется интегралом Фурье,
\begin{equation}
\phi(\mathbf x,t)
=\int\frac{d^3k}{(2\pi)^3}\,
\phi_{\mathbf k}(t)e^{i\mathbf k\cdot\mathbf x}.
\label{eq:c10-fourier-integral}
\end{equation}
Для каждого волнового числа действует один и тот же механизм сдвига к красному и замерзания. \lecturetimestamp{01:46:04--01:46:34}}
\end{classroomqa}

\begin{classroomqa}
\audiencequestion{В конечной Вселенной были бы только дискретные моды?}
\lecturerresponse{Да. Граничные условия в конечном пространственном объёме квантуют допустимые волновые числа. Если объём намного больше обсуждаемых масштабов, шаг между ними может быть настолько мал, что сумма хорошо аппроксимируется интегралом в уравнении~\eqref{eq:c10-fourier-integral}. Непрерывная картина — это идеализация, а не иной локальный механизм. \lecturetimestamp{01:46:34--01:47:01}}
\end{classroomqa}

Теперь сосуществуют два движения. Пространственно усреднённое поле \(\bar\phi(t)\) медленно и классически скатывается вниз по потенциалу. На него накладывается малая квантовая флуктуация,
\[
\phi(\mathbf x,t)=\bar\phi(t)+\delta\phi(\mathbf x,t),
\]
собранная из замороженных мод. Энергия самой флуктуации слишком мала, чтобы давать наблюдаемый контраст вещества. Её важность в том, что она возмущает момент окончания инфляции.

Полный квантовый расчёт для однополевой модели уточняет этот осцилляторный аргумент. Лёгкое скалярное поле получает характерную флуктуацию на каждый логарифмический интервал волнового числа порядка
\begin{equation}
\delta\phi_{\rm exit}\sim\frac{H}{2\pi}.
\label{eq:c10-field-amplitude}
\end{equation}
Вывод нормировки требует квантования мод поля на инфляционном фоне. Однако осцилляторная картина уже дала существенный механизм: вакуумные флуктуации существуют даже в основном состоянии, а расширение превращает их в замороженную картину, выглядящую классической.

\section{Обрыв превращает шум поля в шум плотности}

Сначала пренебрежём \(\delta\phi\). Тогда в любой фиксированный момент поле имеет повсюду одно и то же значение, что на доске изображается вертикальной линией, если пространство \(x\) откладывать вверх, а значение поля \(\phi\) — по горизонтали. По мере течения времени эта линия смещается к обрыву. Вся Вселенная одновременно достигает обрыва, потенциальная энергия инфлатона превращается в частицы, излучение и тепло, экспоненциальное расширение заканчивается, и образовавшаяся плотность энергии начинает разбавляться при обычном расширении.

\begin{figure}[htbp]
\centering
\includegraphics[width=0.88\textwidth]{lecture_10_figure_05.png}
\caption{Схема на доске в координатах «пространство—поле». По вертикальной оси отложено \(x\), по горизонтали — \(\phi\), бледные почти вертикальные профили показывают последовательные конфигурации поля, а заштрихованная граница справа — это обрыв, на котором заканчивается инфляция.}
\label{fig:c10-cliff-board}
\lectureframe{01:50:54.502}
\end{figure}

\begin{figure}[htbp]
\centering
\begin{coursefitfigure}
\begin{tikzpicture}[x=1cm,y=0.82cm,>=Latex]
  \draw[->] (0,0) -- (9.2,0) node[right] {\(\phi\)};
  \draw[->] (0,0) -- (0,5.2) node[above] {space \(x\)};
  \draw[thick,gray!60] (8.2,0.15) rectangle (8.55,4.85);
  \foreach \y in {0.35,0.85,...,4.55} {
    \draw[gray!60] (8.2,\y) -- (8.55,{\y+0.25});
  }
  \node[anchor=south,rotate=90] at (8.38,2.5) {cliff};
  \draw[thick,densely dashed,gray!70] (3.2,0.2) -- (3.2,4.75);
  \draw[thick,blue!45!black,smooth]
    plot coordinates {(5.0,0.2) (5.12,0.75) (4.9,1.25) (5.18,1.8) (4.94,2.35) (5.2,2.9) (4.88,3.45) (5.12,4.05) (5.0,4.75)};
  \draw[very thick,red!65!black,smooth]
    plot coordinates {(7.15,0.2) (7.32,0.7) (7.0,1.2) (7.42,1.75) (7.1,2.3) (7.5,2.85) (7.02,3.35) (7.35,4.0) (7.16,4.75)};
  \draw[->,thick] (3.55,2.5) -- (4.55,2.5);
  \draw[->,thick] (5.55,2.5) -- (6.65,2.5);
  \node[align=center,anchor=south,font=\small] at (3.2,4.8) {uniform\\profile};
  \node[align=center,anchor=south,font=\small] at (5.0,4.8) {frozen\\ripple};
  \node[align=center,anchor=south,font=\small] at (7.2,4.8) {near\\reheating};
\end{tikzpicture}
\end{coursefitfigure}
\caption{Упрощённая реконструкция картины в координатах «пространство—поле». Среднее поле движется к обрыву, тогда как замороженная пространственная рябь сохраняется, поэтому разные точки достигают конца инфляции в немного разные моменты времени.}
\label{fig:c10-cliff-reconstruction}
\end{figure}

Теперь восстановим замороженную вариацию. Линия теперь слегка волнистая: соседние точки имеют немного разные значения \(\phi\). Представьте себе шеренгу людей, пытающихся идти плечом к плечу к обрыву. Некоторые чуть впереди, некоторые чуть позади. Те, кто впереди, первыми переваливаются через край.

Именно этот временной сдвиг и является реальным источником контраста плотности. Область, которая первой переваливается через край, первой и разогревается. Её потенциальная энергия переходит в частицы и тепло, а затем эта энергия успевает разбавиться, пока Вселенная расширяется в послеинфляционную эпоху. Близкая область, которая переваливается через край немного позже, разогревается позже, так что при сравнении этих двух областей более поздняя область оказывается плотнее просто потому, что у неё было меньше времени на разбавление. Поэтому первый набор таких переходов оставляет после себя относительно плотные сгустки. Пока эти сгустки расширяются, рябь с меньшей длиной волны переводит через край ещё один набор точек и создаёт более мелкие сгустки внутри более крупного рисунка. Повторение этого по всей иерархии замороженных мод создаёт структуру внутри структуры. Таким образом, крошечная флуктуация поля важна не из-за собственной энергии; она важна потому, что сдвигает локальный момент, в который происходит разогрев.

Временное соотношение можно записать схематически. Если однородное поле достигает конечного значения \(\phi_{\rm end}\) со скоростью \(\dot{\bar\phi}\), то малое смещение поля изменяет локальное время окончания на
\begin{equation}
\delta t_{\rm end}
\simeq-\frac{\delta\phi}{\dot{\bar\phi}}.
\label{eq:c10-time-delay}
\end{equation}
Чем медленнее скатывание, тем большей задержке соответствует одно и то же \(\delta\phi\). За это дополнительное время область, закончившая раньше, успевает сильнее расшириться и разредиться. С точностью до факторов переноса порядка единицы возмущение кривизны или плотности поэтому определяется выражением
\begin{equation}
\mathcal R
\sim H\,\delta t_{\rm end}
\sim-\frac{H\,\delta\phi}{\dot{\bar\phi}}.
\label{eq:c10-curvature-estimate}
\end{equation}
В сочетании с уравнением~\eqref{eq:c10-field-amplitude} это даёт знакомый масштаб медленного скатывания
\begin{equation}
\mathcal P_{\mathcal R}^{1/2}
\sim\frac{H^2}{2\pi|\dot{\bar\phi}|},
\label{eq:c10-slow-roll-amplitude}
\end{equation}
где для точных соглашений и последующего перехода к \(\delta\rho/\rho\) требуется более полный расчёт возмущений. Чем медленнее скатывается поле, тем больше времени проходит между тем, как соседние области достигают края обрыва, и тем сильнее разрежение между этими моментами. Поэтому одна и та же флуктуация поля оставляет больший отпечаток в плотности, когда \(\lvert\dot{\bar\phi}\rvert\) меньше.

Многие моды пересекают масштаб Хаббла в разные моменты, поэтому многие области разных размеров переразогреваются в слегка разные времена. Возникающее поле плотности содержит переуплотнённые и недоуплотнённые области, вложенные друг в друга на разных масштабах. Его статистика зависит от высоты, наклона и кривизны инфляционного потенциала. Подходящие модели медленного скатывания естественным образом дают почти масштабно-инвариантную амплитуду порядка
\[
\frac{\delta\rho}{\rho}\sim10^{-5},
\]
а детальные спектры CMB позволяют проверять небольшое число космологических параметров на множестве угловых масштабов.

\section{От первичных семян обратно к небу}

Теперь мы можем вернуться к наблюдателю в центре набора оболочек прошлого светового конуса. Сфера последнего рассеяния пересекает трёхмерное поле первичной плотности. В направлениях, где эта оболочка проходит через разные гравитационные потенциалы, температура микроволнового излучения немного различается. WMAP показывает двумерную картину на этой сфере. Эта карта — не фотография нулевого момента времени и не проекция всех структур внутри сферы; это температурная структура на той конкретной оболочке, с которой микроволновые фотоны рассеялись в последний раз.

Внутри этой оболочки те же первичные возмущения продолжают эволюционировать. Гравитация усиливает переуплотнения, каустикоподобные потоки заостряют стенки и филаменты, газ скапливается, и формируются галактики. Таким образом, картина CMB и более поздняя космическая паутина — это две стадии одной истории: замороженные инфляционные флуктуации сначала проявляются как крошечные семена плотности, а позже становятся видимой структурой.

\begin{classroomqa}
\audiencequestion{Какую эпоху представляет наблюдаемый нами сегодня микроволновый фон?}
\lecturerresponse{Фотоны летели примерно \(13\) миллиардов лет, но испущены они были через несколько сотен тысяч лет после горячего Большого взрыва, когда электроны и протоны объединились в нейтральный водород и Вселенная стала прозрачной. Инфляция закончилась несравнимо раньше. Переразогрев создал горячее вещество и излучение; затем Вселенная расширялась и остывала вплоть до последнего рассеяния. Для этого рассуждения надёжно именно утверждение о порядке величины---несколько сотен тысяч лет---, а не точная дата рекомбинации. \lecturetimestamp{02:00:38--02:02:19}}
\end{classroomqa}

\section{Одно небо и космическая дисперсия}

\begin{classroomqa}
\audiencequestion{Предсказывает ли инфляция сообщавшееся выделенное направление или асимметрию север--юг в данных WMAP?}
\lecturerresponse{Из простой инфляционной картины, изложенной здесь, никакое такое выделенное направление не следует. Статус сообщавшихся аномалий на больших угловых масштабах оставался неопределённым, компоненты переднего плана и статистические флуктуации были реальными источниками беспокойства, и для добросовестной оценки требовалась специальная литература, а не сенсационное резюме. Инфляция предсказывает статистические свойства ансамбля; решить, является ли одна необычно выглядящая особенность неба новой физикой, гораздо труднее. \lecturetimestamp{02:02:19--02:03:26}}
\end{classroomqa}

Мелкие угловые особенности встречаются по всему небу много раз. Они дают много выборок, поэтому их распределение и изотропность можно проверять статистически. С особенностью размера квадруполя дело обстоит иначе: на наблюдаемой сфере умещается лишь несколько независимых крупноугловых мод. Поэтому их измеренная мощность может заметно отличаться от среднего по ансамблю даже тогда, когда лежащая в основе теория верна. Эта неустранимая неопределённость выборки называется \emph{космической дисперсией}.

\begin{classroomqa}
\audiencequestion{А как насчёт кажущегося выравнивания нескольких низких мультиполей?}
\lecturerresponse{Это может быть намёком, артефактом переднего плана или анализа, либо маловероятной, но допустимой реализацией. Трудность в том, что мы не можем извлечь ещё одну наблюдаемую Вселенную из того же статистического ансамбля. Монета, давшая 4 орла в 6 бросках, не обязана быть смещённой; это решили бы повторные серии, но космология на самых больших масштабах предоставляет лишь одну такую серию. Аномалия, возможно, заслуживает изучения, однако одно небо само по себе не может дать повторных испытаний, необходимых для решающей проверки статистической значимости. \lecturetimestamp{02:05:31--02:07:36}}
\end{classroomqa}

\section{От квантового движения к космической структуре}

Три универсальные константы определяют планковские масштабы и показывают, насколько далеко привычная физика частиц находится от квантовой гравитации. Микроволновое небо даёт точный космологический ориентир: почти изотропность, амплитуду флуктуаций порядка \(10^{-5}\) и структуру на многих масштабах. Квантовая механика гарантирует, что у каждой моды поля есть флуктуации нулевой точки. Инфляция растягивает каждую моду инфлатона, пока её физическая длина волны не станет порядка \(H^{-1}\), после чего хаббловское демпфирование подавляет её осцилляцию и оставляет эффективно замороженную амплитуду.

Однородный инфлатон продолжает медленно скатываться на фоне этого замороженного шума. Когда поле достигает конца своего плато, соседние области не переразогреваются в точности одновременно. Раньше закончившие области разрежаются дольше; закончившие позже остаются более плотными. Это превращает \(\delta\phi\) в первичную картину плотности и кривизны. Затем гравитация усиливает эту картину, и из неё возникают галактики и космическая паутина, тогда как оболочка последнего рассеяния сохраняет своё более раннее состояние на микроволновом небе. Следовательно, микроскопическое движение нулевой точки, космическое расширение и несинхронный переразогрев — не отдельные истории. Вместе они объясняют, как инфляционная Вселенная может быть одновременно необычайно гладкой и достаточно структурированной, чтобы мы могли существовать.